% ==================== 七、模型评价与推广 ====================
\section{模型评价与推广}

\subsection{模型优点}

\begin{enumerate}[itemsep=4pt,topsep=4pt]
  \item \textbf{统一而递进的建模框架。}
        四问共享"预测 $\to$ 场景 $\to$ 线性规划 $\to$ 滚动执行"的基本结构，
        并逐层增加不确定性维度（电价已知 $\to$ 负载/光伏未知 $\to$ 多时点信息 $\to$ 电价亦未知）。
        这种设计使得各问之间的差异\textbf{可归因、可对照}，
        避免了"四问四套模型、无法横向比较"的常见问题。

  \item \textbf{报童结构的识别提供了经济学解释。}
        我们把问题二识别为报童模型，导出临界分位 $F^*=0.80$，
        说明"计划量应取 80\% 分位而非均值"。
        这使参数设定有理论依据，而非依赖盲目试错；
        也解释了为何早期"按点预测制定计划"的方案会产生大量昂贵的紧急购电。

  \item \textbf{LP 松弛紧性的严格证明（问题一）。}
        用等 SOC 扰动证明：在往返效率 $0.81<1$ 且该时段有购电的前提下，
        同时充放电必劣于拆开。由此\textbf{无需二元变量}，
        问题规模保持在线性规划量级，而 LP 与 MILP 数值结果完全一致
        （$35\,126.95$ 元，差 $0.000000$），证明了推导的正确性。

  \item \textbf{非前视由变量结构保证。}
        问题三的场景树中，同一节点内的所有场景\textbf{共享同一组控制变量}，
        使模型在结构上不可能"针对已知场景分别决策"。
        这比仅添加额外约束更为稳健，也便于审阅者验证。

  \item \textbf{电价—净负载联合场景的正确性论证。}
        我们指出 $\E[p\cdot(N-z)_+]\neq\E[p]\E[(N-z)_+]$，
        并据此采用同历史日配对的场景构造，保留了"高电价与缺电同时发生"的风险耦合。

  \item \textbf{严格的因果性与可复现性。}
        所有预测与融合权重均只使用决策时刻之前的数据；
        程序对 $48\,096$ 个时段逐段核验能量平衡、效率递推、SOC 边界、功率上限、
        互斥性、跨日连续、初末状态、缺口与费用，并反读 Excel 输出。
    \item \textbf{诚实的对照报告。}
        所有未胜出的方案（半衰期加权、短窗口、在线选窗、平均电价模型、
        多源信息、$k$ 参数扫描等）均在结果与附录中保留，
        便于审阅者评估结论的稳健性。
\end{enumerate}

\subsection{模型局限}

\begin{enumerate}[itemsep=4pt,topsep=4pt]
  \item \textbf{非全局最优}：本文为有限经验场景上的线性规划加滚动执行，
        不宣称全年真实随机控制问题的全局最优解。
        小型二叉场景树、次日确定性近似、反复滚动重算都是对真实问题的近似。

  \item \textbf{结果均为同年回测}：所有费用数字来自 2025 年 2--12 月的滚动回测，
        不证明在其他年份仍最优。例如问题三中"$18$ 点更新无收益"是本年度数据结论。

  \item \textbf{预测内核参数的选型过程未完全审计}：
        负载的"最近 4 次同星期 $+$ 近 5 日漂移"、光伏的"近 7 日 $+$ 近 28 日漂移"
        等参数继承自前期研发，
        本次因果性检查只能验证"运行时未使用未来数据"，
        不能证明当初的参数选择未使用过测试年份信息。

  \item \textbf{题面若干处的建模解释需公开声明}：
        问题三"取消部分退原价、另付 50\% 违约费"是本文采用的结算解释；
        若改为"不退原价、另罚 50\%"，则同一调度方案的费用会变化，
        且应重新优化而非仅调整报表公式。该敏感性已写入核验文件。

  \item \textbf{储能模型简化}：未计入折旧、寿命衰减、温度与功率相关性，
        也未考虑电池的最小充放电持续时间等工程约束。

  \item \textbf{场景样本量受限于历史长度}：2 月 1 日仅 24 条残差样本，
        场景分布的尾部刻画精度有限。
\end{enumerate}

\subsection{模型推广}

\begin{enumerate}[itemsep=3pt,topsep=3pt]
  \item \textbf{引入外部气象数据}。当前光伏预报误差是主要成本来源
        （完美预见下界与现实的差距约 $179$ 万元）；
        引入辐照度、云量、温度等特征，或按天气类型分别建模，
        有望显著降低预报误差。

  \item \textbf{自适应场景树深度}。当前固定 4 层、15 节点，
        可依据样本量与预报修订的方差结构自适应确定分支深度，
        在样本充足时提供更细的信息分辨率。

  \item \textbf{电价预测的增强}。可引入电价曲线的日内形状聚类、
        天气耦合与节假日效应，改善 4-2/4-3 中电价预测的精度
        （该误差对应约 $1.0\%$ 的费用）。

  \item \textbf{储能全寿命周期成本建模}。加入循环次数与放电深度的折旧模型，
        可避免"为套利而过度循环"的短视调度。

  \item \textbf{多微网协同}。若扩展至多微网共享储能或互联互济，
        本文的场景树框架可直接推广为多主体随机博弈结构。
\end{enumerate}
